\begin{frame}{TD vs MC: 편향과 분산}
  \small
  \begin{tabular}{@{}lll@{}}
    \toprule
     & 몬테카를로 & 시간차 학습(TD) \\
    \midrule
    목표값 & 실제 리턴 $G_t$(끝까지) & $r+\gamma V(s')$(한 스텝+추정치) \\
    갱신 시점 & 에피소드 종료 후 & 매 스텝 즉시 \\
    편향/분산 & 편향 없음(불편), \textbf{분산 큼}
               & 편향 있음(추정치 사용), \textbf{분산 작음} \\
    \bottomrule
  \end{tabular}
  \vskip1em
  \begin{itemize}
    \item MC의 $G_t$는 실제로 관찰된 보상의 합이라 편향이 없으나,
      에피소드 하나가 \textbf{긴 시퀀스의 모든 확률성}을 통째로 담고 있어
      분산이 크다.
    \item TD의 목표값은 $V(s')$이 \textbf{아직 부정확한 추정치}라 편향이
      있으나, 한 스텝의 무작위성만 반영해 분산이 훨씬 작다.
  \end{itemize}
  \vskip0.4em
  {\footnotesize Week 4.1의 편향-분산 트레이드오프가 여기서도 다른 형태로
  다시 등장한다.}
\end{frame}
